%%%%%%%%%%%%%%%%%%%%%%%%%%%%%%%%%%%%%%%%%
% University Assignment Title Page
% LaTeX Template
% Version 1.0 (27/12/12)
%
% This template has been downloaded from:
% http://www.LaTeXTemplates.com
%
% Original author:
% WikiBooks (http://en.wikibooks.org/wiki/LaTeX/Title_Creation)
%
% License:
% CC BY-NC-SA 3.0 (http://creativecommons.org/licenses/by-nc-sa/3.0/)
%
% Instructions for using this template:
% This title page is capable of being compiled as is. This is not useful for
% including it in another document. To do this, you have two options:
%
% 1) Copy/paste everything between \begin{document} and \end{document}
% starting at \begin{titlepage} and paste this into another LaTeX file where you
% want your title page.
% OR
% 2) Remove everything outside the \begin{titlepage} and \end{titlepage} and
% move this file to the same directory as the LaTeX file you wish to add it to.
% Then add \input{./title_page_1.tex} to your LaTeX file where you want your
% title page.
%
%%%%%%%%%%%%%%%%%%%%%%%%%%%%%%%%%%%%%%%%%
%\title{Title page with logo}
%----------------------------------------------------------------------------------------
%   PACKAGES AND OTHER DOCUMENT CONFIGURATIONS
%----------------------------------------------------------------------------------------

\documentclass[12pt]{article}
\usepackage[german]{babel}
\usepackage[T1]{fontenc} % for typewriter braces
\usepackage[utf8x]{inputenc}
\usepackage{amsmath}
\usepackage{graphicx}
\usepackage[colorinlistoftodos]{todonotes}
% Simple directory trees
\usepackage{dirtree}
\usepackage{nameref}
\usepackage{listings}

% \meta command
\usepackage{textcomp}
\usepackage{hyperref}
\providecommand\meta[1]{\textlangle{\itshape #1\/}\textrangle}

\begin{document}

    \begin{titlepage}

        \newcommand{\HRule}{\rule{\linewidth}{0.5mm}} % Defines a new command for the horizontal lines, change thickness here

        \center % Center everything on page

        % Heading
        \textsc{\LARGE GFOS Innovationsaward 2018}\\[1.5cm]
        \textsc{\Large Helmholtz-Gymnasium Essen}\\[0.5cm]
        \textsc{\large Dokumentation}\\[0.5cm]

        % Main title
        \HRule \\[0.4cm]
        { \huge \bfseries ZNews Server\\ Bedienungsanleitung}\\[0.4cm]
        \HRule \\[1.5cm]

        % Authors
        \begin{minipage}{0.4\textwidth}
            \begin{flushleft}
                \large
                \emph{Autoren:}\\
                Julian Koch\\
                Mats Holtbecker
            \end{flushleft}
        \end{minipage}
        ~
        \begin{minipage}{0.4\textwidth}
            \begin{flushright}
                \large
                \emph{Aufsicht:} \\
                Dennis Gro"skamp\\
                \hfill
                % Empty line
            \end{flushright}
        \end{minipage}\\[2cm]

        % If you don't want a supervisor, uncomment the two lines below and remove the section above
        %\Large \emph{Author:}\\
        %John \textsc{Smith}\\[3cm] % Your name

        % Date
        {\large \today}\\[2cm]

        % GFOS Logo
        \includegraphics[scale=0.7]{../gfos_logo_edited.jpg}\\[1cm] % Include a department/university logo - this will require the graphicx package

        %----------------------------------------------------------------------------------------

        \vfill
        % Fill the rest of the page with whitespace

    \end{titlepage}


    %% We don't need an abstract in this page
    %% We'll add that to Dokumentation.pdf
    %\begin{abstract}
    %    ZNews besteht weitestgehend aus zwei sich gegen"uberstehenden
    %    Komponenten: dem Server und dem Webinterface.
    %    In diesem Dokument wird die Funktionsweise des Servers
    %    dargestellt und erl"autert.
    %\end{abstract}

    \newpage
    \tableofcontents
    \newpage

    \section{Installation}

    ZNews Server wird als archivierte Datei geliefert.
    Nach deren Download kann sie mit einem Archivierungsprogramm
    wie WinRAR oder 7-Zip ausgepackt werden.

    Die Struktur des Hauptordners ist in Abb.~\ref{fig:dirtree1}
    dargestellt.

    \begin{enumerate}
        \item Extrahieren Sie nach deren Download die Datei \emph{ZNews.zip}.
        Die Struktur des Hauptordners ist in Abb.~\ref{fig:dirtree1} dargestellt.
        \item Starten Sie den Server mit dem Startskript\footnote{%
        Windows: \emph{start.bat}, Mac: \emph{start.command}}.
        \item Nach Erscheinen der Nachricht ``Server started'' (o."a.),
        Schreiben Sie ``end'' und dr"ucken Sie \meta{Enter}.
        \item Nach Beenden des Programms sind die in Abb.~\ref{fig:dirtree2}
        gezeigten Dateien und Ordner neu erstellt worden.
        Konfigurieren Sie den Server wie in Abschnitt~\ref{sec:config} ``\nameref{sec:config}'' beschrieben.
        \item Sie k"onnen nun jederzeit den Server mit dem Startskript starten.
    \end{enumerate}

    \begin{figure}[htb]
        \dirtree{%
        .1 /.
        .2 server.jar.
        .2 start.bat.
        .2 start.command.
        .2 docs/.
        .3 de/.
        .4 Bedienungsanleitung.pdf.
        .4 Server Bedienungsanleitung.pdf.
        .4 Dokumentation.pdf.
        }
        \caption{\label{fig:dirtree1}Aufbau des heruntergeladenen Archivs}
    \end{figure}

    \begin{figure}[htb]
        \dirtree{%
        .1 /.
        .2 logs/.
        .2 static\textunderscore web/.
        .2 auth.json.
        .2 config.properties.
        .2 newsletters.json.
        .2 registrations.json.
        }
        \caption{\label{fig:dirtree2}Neu erstellte Dateien und Ordner}
    \end{figure}

    \pagebreak
    \section{Konfiguration}%
    \label{sec:config}

    ZNews Server beinhaltet viele Konfigurationsm"oglichkeiten,
    allesamt in \emph{config.properties}.
    Dieser Abschnitt zeigt die wichtigsten Konfigurationen
    und bietet au"serdem eine vollst"andige Referenz am Ende.

    \subsection{Der Port}

    \begin{align*}
        \emph{Schl"ussel: } & \lstinline{port}\\
        \emph{Standardwert: } & \lstinline{8080}
    \end{align*}
    %% Consider this other, non-centered version:
    %\hspace*{2cm}\emph{Schl"ussel: } \lstinline{port}\\
    %\hspace*{2cm}\emph{Standardwert: } \lstinline{8080}
    %\vspace*{0.2cm}

    ZNews Server beansprucht einen Port, um eingehende
    Anfragen bearbeiten zu k"onnen. Standardm"a"sig ist
    \emph{8080} eingestellt, jedoch ist es empfohlen,
    diesen Wert zu "andern.

    \subsection{Email-Einstellungen}

    ZNews Server muss dazu in der Lage sein, E-Mails
    zu versenden. Dazu muss z.B. ein SMTP-Server eingestellt
    werden, durch welchen die E-Mails geschickt werden.

    Im Folgenden wird ein Gmail-Account mit der
    E-Mail-Adresse \emph{john.smith@gmail.com} und
    dem Passwort \emph{12345} verwendet:

    \begin{lstlisting}[]%@formatter:off
email.protocol=smtp
email.port=465
email.host=smtp.gmail.com
email.auth=true
email.auth.usr=john.smith@gmail.com
email.auth.pw=12345
    \end{lstlisting}%@formatter:on

    F"ur die Verwendung von anderen Servern, konsultieren Sie
    bitte deren spezifischen Anleitungen.

    \subsection{Log-Level}

    \begin{align*}
        \emph{Schl"ussel: } & \lstinline{log.filter}\\
        \emph{Standardwert: } & \lstinline{out}\\
        \emph{M"ogliche Werte: } & \lstinline{dev|debug|out|warn|err|fatal}
    \end{align*}

    Das Log-Level spezifiziert, welche Logging-Ausgaben tats"achlich
    ausgegeben werden. Zum Beispiel werden bei \lstinline{warn}
    nur Warnungen, Fehler (\lstinline{err}) und schwerwiegende Fehler
    (\lstinline{fatal}) ausgegeben.

    \subsection{Lesbare JSON-Ausgabe}

    \begin{align*}
        \emph{Schl"ussel: } & \lstinline{pretty-print-json}\\
        \emph{Standardwert: } & \lstinline{false}
    \end{align*}

    Standardm"a"sig werden die von ZNews ben"otigten Daten
    in den Dateien \emph{auth.json}, \emph{newsletters.json} und \emph{registrations.json}
    gespeichert\footnote{%
    Alternativ kann auch eine MySQL-Datenbank verwendet werden, siehe % // TODO: Add option for MySQL
    }, jedoch sind sie nicht besonders leserlich.

    Die Lesbarkeit wird verbessert, wenn Sie diese Option aktivieren.


    %Comments can be added to the margins of the document using the \todo{Here's a comment in the margin!} todo command, as shown in the example on the right. You can also add inline comments too:
    %
    %\todo[inline, color=green!40]{This is an inline comment.}


\end{document}